\chapter{Evaluation and Conclusion}
% A summary of what the project has achieved. Make sure that you
% address each goal set out in the Introduction chapter, to show that you
% have achieved what you claimed you would. Don’t leave any loose
% ends.
% – A critical evaluation of the results of the project (e.g., how well were
% the goals met, is the application fit for purpose, has good design and
% implementation practice been followed, was the right implementation
% technology chosen and so on).
% – Future work. How could the project be developed if you had another 6
% months. Take care to differentiate between what you have done to
% satisfy your stated project goals, and work that could be done to meet
% extended goals.
% – Wrap-up and final thoughts on your project.
\section{Achievements}
This project has completed all of the goals set out in the introduction providing evidence that this proof-of-concept is worth persuing for future development. The initial comparison between candidate medical LLMs showed Med42-v2 to have the optimal accuracy-speed tradeoff out of the candidate models, demonstrating the effectiveness of OpenVINO's model quantisation framework. A real time transcription-diarisation pipeline was implemented using the \texttt{Whisper-Medium} and \texttt{ECAPA-TDNN} models, connected to the summarisation pipeline. The appraisal entry system was created based on feedback from Dr. Atia Rafiq, and a guidelines retrieval component was realised. An iterative approach to the UI design enabled an appealing and user friendly end product, native to Windows. Finally, the application design using the orchestator and pipeline design patterns, and sub-implementations of the controllers resulted in an extensible application. 

Outside of the goals, I produced 6 demonstration videos, one of which was presented at the shard at a Dell event for Intel and NHS stakeholders. The final build of the application was presented at a an Intel sponsored HP booth at the NHS Strategy Summit on the 29th April 2026. I featured on the UCL Computer Science LinkedIn page while conducting demonstrations to healthcare professionals and presented to 50 healthcare professionals at the UCL Computer Science Showcase 2026. Finally, the application has been published through my supervisor's Motion Input Games platform and is free to download (strictly for demonstration purposes) at \url{https://www.motioninputgames.com/software/AIPC/health/Clinical-Summarisation-AI-PC_1.0.10.0_x64.msix}.

\section{Critical Evaluation of the Results}
The project was successful in demonstrating that an offline clinical documentation workflow can be implemented on Intel AI PC hardware using a combination of quantised language models, speech recognition, speaker diarisation, and local retrieval. The implemented system satisfies the intended proof-of-concept scope as it records a consultation, produces a structured summary, stores supporting data locally, and allows clinicians to attach contextual guidance and appraisal material without relying on cloud infrastructure. 

The application design followed sound software engineering principles, particularly through separation of concerns. The concurrent architecture isolated the user interface from computationally expensive inference tasks, which was important for maintaining responsiveness during a consultation. The data management approach also separated transient workflow data from persistent records such as appraisal JSON files and the SQLite database. Interface abstraction enabled the use of mock objects during testing, which improved the ability to evaluate components in isolation.

Some design decisions also introduced limitations. The choice to keep user interface state within a single \texttt{MainWindow} class, even when split across multiple translation units, was understandable for a sequential proof-of-concept but is less maintainable than a full MVVM architecture if the interface grows further. The evaluation is also stronger on functional demonstration than on formal quantitative validation. For example, this project presents evidence that the transcription and summarisation pipeline works in practice, but there is limited end-to-end measurement of consultation time saved, documentation quality, or retrieval relevance across a large dataset. Formal user acceptance testing and a clinical trial would be required to determine whether the application provides measurable benefit to GPs compared to other solutions. Furthermore, additional research into alternative models such as Whisper-Turbo or a mechanism to enable user uploaded models would be required to establish the correctness of the models chosen. As result, this project demonstrates feasibility more than clinical time saved.

The development process was generally effective, particularly where regular meetings with supervisors and clinicians helped keep the work aligned with practical requirements. However, an early assumption that NICE API access would be feasible proved incorrect. This led to redesign work later in the project and resulted in a less integrated retrieval approach based on clinician uploaded guidance documents. 

\section{Future Work}
Although the proof-of-concept met its core requirements, several important extensions remain out of scope. The most significant short-term goal would be a structured clinical evaluation. A planned trial with Dr Tazeen Ahmed at St. George's Hospital was blocked because of the March-April 2026 cyber attack, but a future study of this kind would be valuable. Measuring time saved, documentation burden, usability, and perceived accuracy during real consultations would enable a comparative study against current solutions. It would also identify where the current workflow creates friction for clinicians.

If access to the NICE API became available, a hybrid online-offline architecture could download updated guidance at controlled intervals while preserving local inference during use. This would enable more rigorous retrieval evaluation over a larger and better defined document collection.

Further feature development could also build on the existing pipeline structure. Translation was discussed during the project but remained out of scope however, adding a translation stage could improve accessibility for multilingual consultations. Additional recording functionality, such as hands-free activation, could better match how clinicians work in practice. Speech-to-text input for appraisal notes would also be a natural extension of the current transcription pipeline as suggested by Dr. Atia Rafiq.

Finally, future hardware and OpenVINO improvements may allow larger local models and longer prompts, which could improve summarisation quality. The direction of current AI PC development and potential quantisation improvements could allow larger summarisation models to be tested. Future summarisation work could also include patient-friendly explanations of diagnoses and tighter integration of retrieved guidance into the generated summary.

\section{Conclusion}
This project demonstrated that an offline, Windows-native clinical summarisation system is feasible. The final application combines local speech transcription, diarisation, summarisation, appraisal support, and retrieval of medical guidance into a coherent proof-of-concept workflow. While the evaluation does not yet establish clinical effectiveness at deployment scale, it does show that the underlying architecture and implementation are strong enough to support further development.

Personally, I learned a new programming language, development environment, packaging environment and font-end stack. This lead to lots design and planning, with many teething problems throughout the highly challenging yet enjoyable development journey. Overall, this project acts as the first building block towards offline AVT solutions for Intel and the NHS to reduce the time taken and fatigue caused by note taking and documentation within clinical consultations throughout healthcare.